\begin{textsample}{POS Dim 1 – rr – Score 40.00 – rr/rr\_ef\_63.txt}  \label{ex:f1_pos_218}
5.2 METODOLOGIAS E ESTRATÉGIAS DIDÁTICO PEDAGÓGICAS O Documento Curricular de Roraima - DCR, construído em regime de colaboração entre estado e municípios, traz a responsabilidade de representar o território em sua diversidade de saberes e vivências culturais.
Na Base Nacional Comum Curricular-BNCC, as \textbf{competências} são definidas como conjunto de conceitos e procedimentos, desdobradas em práticas cognitivas e \textbf{socioemocionais}, no tratamento didático proposto para as etapas da Educação Infantil e Ensino Fundamental, da Educação Básica.
Cumpre destacar que as discussões acerca do desenvolvimento educacional dos \textbf{alunos} do estado requerem atenção criteriosa quanto às metas, estratégias e ações traçadas junto aos sistemas de ensino, as equipes pedagógicas e professores, quanto aos pressupostos didáticos e \textbf{metodológicos} \textbf{inerentes} a prática docente no contexto da sala de \textbf{aula}.
Dessa forma, refletir a prática pedagógica na perspectiva do DCR, requer estudos e discussões das diretrizes institucionais e Projetos Políticos Pedagógicos-PPP de todas as redes de ensino, a partir do alinhamento e coerência com os objetivos de aprendizagem e das habilidades, por meio das \textbf{Competências} Gerais e Específicas da BNCC.
Estas mudanças ampliam as perspectivas de pensar a formação continuada dos professores, fundamentada nos processos e procedimentos \textbf{inerentes} a realidade local e ao papel docente na sala de \textbf{aula}.
Ao pensar os princípios norteadores da ação docente, precisamos buscar e discutir meios de desenvolvimento de esforços pedagógicos comuns, provocando situações que promovam sujeitos \textbf{aprendizes} em sua atividade cotidiana, a partir de metodologias adequadas às necessidades e interesse dos \textbf{alunos}, rediscutindo a identidade e diversidade sociocultural da população escolar, organizadas nas práticas didáticas dos currículos escolares.
Neste contexto, apesar de vários documentos anteriores a BNCC, como por exemplo, as DCNs, chamarem a atenção para necessidade de superarmos no trabalho escolar as ideias de disciplinaridade, por meio de um trabalho mais integrado, contextualizado e interdisciplinar, a prática da escola, tem se modificado de forma muito lenta.
Segundo Ocampo, Santos e Fomer ( 2016, p. ) “ A \textbf{interdisciplinaridade} possibilita não apenas a interação de conteúdo, mas também a interação entre pessoas, já que essa perspectiva tem \textbf{potencial} para motivar outros professores que compõem o corpo docente ”.
Neste sentido \textbf{exige} que os participantes partilhem conhecimentos e experiências, adotando um comportamento de diálogo e respeito.
Esse partilhamento é essencial e \textbf{central} na ideia de protagonismo do \textbf{aluno} e professor, como também da empatia, destacado como uma das \textbf{competências} gerais da BNCC.
Ainda ressaltando a importância e a inserção da \textbf{interdisciplinaridade} no trabalho educacional, Moraes et all. ( 2014 ) afirmam que essa perspectiva está relacionada “ a uma mudança de paradigma no âmbito do currículo que vem incorporando temas como caos, complexidade, multiculturalismo, dialogicidade, intersubjetividade ”.
Como também contrapõe a linearidade e fragmentação características de uma visão tradicional que vêm cedendo espaço à integração, \textbf{contextualização} e busca da \textbf{totalidade} no conhecimento. \textbf{Fomentar} essa \textbf{totalidade} do conhecimento é cada vez mais necessário no mundo de mudanças, em que as respostas \textbf{simples} e disciplinares não dão conta de responder e explicar adequadamente a realidade \textbf{atual}.
Os \textbf{autores} citados, destacam também que a \textbf{interdisciplinaridade} é “ [... ] ao mesmo tempo uma \textbf{abordagem} epistemológica e \textbf{metodológica} desde que não só desafia as fronteiras disciplinares como também acrescenta novas \textbf{análises} da realidade ” ( MORAES et. all, 2014 ).
Precisamos então, buscar \textbf{caminhos} para reestruturar a forma de trabalho da escola, privilegiando um trabalho mais coletivo e menos individual.
É fundamental buscar metodologias que agucem seus diversos \textbf{atores} a uma participação mais ativa e compartilhada, superando a especialização que impregna o trabalho escolar.
Dentre essas metodologias poderíamos citar : os projetos educativos e o estudo por problemáticas.
Porém, não temos o \textbf{intuito} de direcionar ou limitar os procedimentos que devam ser usados pelo professor e pela escola.
O mais importante é que essas decisões \textbf{metodológicas} sejam feitas no coletivo que encontrará a melhor forma de desenvolver seu trabalho, de acordo com sua realidade, e que contemple essas reflexões nos PPPs das escolas, e consequentemente na formação continuada dos professores.
Essas decisões deverão ter como base as \textbf{competências} gerais da BNCC, como por exemplo, a de : > Exercitar a curiosidade intelectual e recorrer à \textbf{abordagem} própria da \textbf{ciência}, incluindo a investigação, a reflexão, a \textbf{análise} \textbf{crítica}, a imaginação e a criatividade, para investigar causas, elaborar e testar hipóteses, formular e resolver \textbf{problemas} e cria soluções ( inclusive tecnológicas ) com base nos conhecimentos de diferentes áreas ( BRASIL, 2013, p. 9 ).
Corroborando com essa ideia, Vieira et all ( 2018 ) ressaltam que a atuação do professor como \textbf{pesquisador}, \textbf{tanto} colabora para sua auto formação como contribui para sua prática docente.
Além disso, essa forma de trabalho é propicia para aproveitar a \textbf{capacidade} de aprendizagem dos \textbf{alunos}, como também \textbf{despertar} o interesse destes pelos objetos de conhecimentos.
Considerando que as escolas se constituem em ambientes repletos de vivências específicas, decorrentes da diversidade sociocultural que circula nesse espaço, torna -se fundamental que na educação escolar sejam elaborados PPPs viáveis, com enfoques alternativos, considerando a realidade de quem ensina e de quem aprende, estruturando -se nos cenários onde se contextualizam saberes e vivências específicas, se distanciando do ensino e aprendizagem tradicionalmente instituídos nos currículos e difundidos no dia a dia da sala da \textbf{aula}.
Nesse sentido, a educação escolar deve se \textbf{preocupar} em oportunizar a igualdade educacional, respeitar a diversidade e promover a equidade, e para que isso ocorra, as Diretrizes Curriculares Nacionais – DCNs ( BRASIL, 2013 ), orientam que nas instituições escolares aconteçam a inclusão social, garantindo que a diversidade humana, social, cultural, \textbf{econômica} dos grupos historicamente excluídos seja respeitada.
Trata -se de observar : > as questões de classe, gênero, raça, etnia, geração, constituídas por \textbf{categorias} que se entrelaçam na vida social – pobres, mulheres, afrodescentendes, indígenas, pessoas com deficiência, as populações do campo, os de diferentes orientações sexuais, os sujeitos albergados, aqueles em situação de rua, em privação de liberdade – todos que compõem a diversidade que é a sociedade brasileira e que começam a ser contemplados pelas políticas públicas ( BRASIL, 2013, p. 16 ).
Desse modo, é fundamental que as propostas educacionais desenvolvidas garantam a aprendizagem de conhecimentos e \textbf{competências} necessárias para que os estudantes se posicionem com sensatez frente às situações do cotidiano.
Assim, as aprendizagens devem assegurar o domínio nas diferentes áreas do conhecimento e também a formação humana integral, objetivando a construção de uma sociedade mais justa, democrática e inclusiva.
Nessa perspectiva, torna -se importante estimular o protagonismo dos \textbf{alunos} no seu processo de ensino e aprendizagem, favorecendo o desenvolvimento de \textbf{competências} e habilidades, no \textbf{intuito} de que possa relacionar o saber construído com o conhecimento posto em uso, indo além das abstrações desconexas.
Neste sentido, é fundamental romper com a tradição formalista dos procedimentos de ensino, isso não significa que os conhecimentos formais devam ser ignorados, mas é importante buscar outros meios de ensino, dando \textbf{voz} ao estudante.
Assim sendo, o PPP deve considerar o atendimento às singularidades dos estudantes, considerando as suas especificidades \textbf{tanto} históricas como familiar, reconhecendo que as suas necessidades são diferentes.
Por \textbf{conseguinte}, para atender as necessidades educacionais de uma sociedade diversa é necessário refletir sobre os aspectos que envolvem a vida do estudante : seu povo, sua história, cultura, tradições, saberes, bem como expectativas e necessidades locais.
Nesse contexto, as ações educacionais desenvolvidas nas escolas devem considerar “ [... ] as necessidades, as possibilidades e os interesses dos estudantes, assim como suas identidades linguísticas, étnicas e culturais ” ( BRASIL, 2017, p. 15 ).
Em \textbf{vista} disso, a dinâmica social e cultural existente no estado de Roraima \textbf{exige} que as metodologias e estratégias didático-pedagógicas sejam diversificadas, como forma de garantir a valorização da diversidade que \textbf{permeia} a nossa sociedade, uma vez que ambas, sociedade e escola, convivem com culturas locais, regionais, nacionais e internacionais.
Portanto, é \textbf{relevante} que nas intervenções educacionais sejam desenvolvidas estratégias e procedimentos de ensino diferenciados, os quais valorizem a diversidade sociocultural dos estudantes e promovam a sua autonomia, para que possam ser motivados e engajados no fortalecimento das suas aprendizagens.
Entre as ações de intervenção que poderão \textbf{emergir} desse processo, considera -se a organização e implementação de projetos, pesquisas e \textbf{sequências} didáticas que possam orientar o trabalho docente, considerando a possibilidade de aliar \textbf{teoria} e prática, valorizar os conhecimentos e a diversidade sociocultural e dar significado aos objetos de conhecimentos organizados no currículo escolar da Educação Infantil e Ensino Fundamental da Educação Básica do estado.

% matched lemmas: abordagem, aluno, análise, aprendiz, ator, atual, aula, autor, caminho, capacidade, categoria, central, ciência, competência, conseguinte, contextualização, crítico, despertar, econômico, emergir, exigir, fomentar, inerente, interdisciplinaridade, intuito, metodológico, permear, pesquisador, potencial, preocupar, problema, relevante, sequência, simples, socioemocional, tanto, teoria, totalidade, vista, voz
\end{textsample}
